% GENERATED FILE. Edit canonical lessons, then run npm run generate:books.
% canonical-source-hash: fnv1a64:6869051f4e532f65
% canonical-lessons: SA-C95-madhye, SA-C95-parsve, SA-C95-sarvatra, SA-C95-anyatra, SA-C95-vithi

\chapter{In the middle, Beside, Everywhere, Elsewhere, Street}
\label{ch:sa-in-the-middle-beside-everywhere-elsewhere-street}
\hlchaptermodality{95}

\begin{chapteropening}
\textbf{By the end of this chapter:} \emph{I can say in the middle, beside, everywhere, elsewhere and a street.}

Complete the last lesson of chapter 95: I can say in the middle, beside, everywhere, elsewhere and a street.
\end{chapteropening}

\section[\texorpdfstring{\sk{मध्ये} (madhye)}{मध्ये (madhye)}]{\sk{मध्ये} (madhye) — in the middle}

\label{lesson:SA-C95-madhye}



\begin{quote}
Before the new one: say the Sanskrit for in front, then the Sanskrit for behind.
\end{quote}

\subsection*{You'll want to know: \sk{मध्ये}}
\textbf{\sk{मध्ये}} — \emph{madhye} — \textquotedblleft{}in the middle\textquotedblright{}.

\begin{grammarlens}[title={What you've built}]
One more word for this chapter.
\end{grammarlens}

\subsection*{Your turn}
\begin{itemize}
\raggedright
  \item \emph{Say it:} \emph{madhye}
  \item \emph{Say it:} \emph{madhye}, once more
  \item \emph{Say it:} say \emph{madhye}
  \item \emph{Recall:} say the Sanskrit for in front, then the Sanskrit for behind, then say \emph{madhye} again
\end{itemize}

\begin{tcolorbox}[breakable,colback=teal!4,colframe=teal!35!black,title={Before you move on}]
What does \sk{मध्ये} mean? (\textquotedblleft{}In the middle\textquotedblright{}.) Say it once more.
\end{tcolorbox}
\section[\texorpdfstring{\sk{पार्श्वे} (pārśve)}{पार्श्वे (pārśve)}]{\sk{पार्श्वे} (pārśve) — beside}

\label{lesson:SA-C95-parsve}



\begin{quote}
Before the new one: say the Sanskrit for behind, then the Sanskrit for in the middle.
\end{quote}

\subsection*{You'll want to know: \sk{पार्श्वे}}
\textbf{\sk{पार्श्वे}} — \emph{pārśve} — \textquotedblleft{}beside\textquotedblright{}.

\begin{grammarlens}[title={What you've built}]
One more word for this chapter.
\end{grammarlens}

\subsection*{Your turn}
\begin{itemize}
\raggedright
  \item \emph{Say it:} \emph{pārśve}
  \item \emph{Say it:} \emph{pārśve}, once more
  \item \emph{Say it:} say \emph{pārśve}
  \item \emph{Recall:} say the Sanskrit for behind, then the Sanskrit for in the middle, then say \emph{pārśve} again
\end{itemize}

\begin{tcolorbox}[breakable,colback=teal!4,colframe=teal!35!black,title={Before you move on}]
What does \sk{पार्श्वे} mean? (\textquotedblleft{}Beside\textquotedblright{}.) Say it once more.
\end{tcolorbox}
\section[\texorpdfstring{\sk{सर्वत्र} (sarvatra)}{सर्वत्र (sarvatra)}]{\sk{सर्वत्र} (sarvatra) — everywhere}

\label{lesson:SA-C95-sarvatra}



\begin{quote}
Before the new one: say the Sanskrit for in the middle, then the Sanskrit for beside.
\end{quote}

\subsection*{You'll want to know: \sk{सर्वत्र}}
\textbf{\sk{सर्वत्र}} — \emph{sarvatra} — \textquotedblleft{}everywhere\textquotedblright{}.

\begin{grammarlens}[title={What you've built}]
One more word for this chapter.
\end{grammarlens}

\subsection*{Your turn}
\begin{itemize}
\raggedright
  \item \emph{Say it:} \emph{sarvatra}
  \item \emph{Say it:} \emph{sarvatra}, once more
  \item \emph{Say it:} say \emph{sarvatra}
  \item \emph{Recall:} say the Sanskrit for in the middle, then the Sanskrit for beside, then say \emph{sarvatra} again
\end{itemize}

\begin{tcolorbox}[breakable,colback=teal!4,colframe=teal!35!black,title={Before you move on}]
What does \sk{सर्वत्र} mean? (\textquotedblleft{}Everywhere\textquotedblright{}.) Say it once more.
\end{tcolorbox}
\section[\texorpdfstring{\sk{अन्यत्र} (anyatra)}{अन्यत्र (anyatra)}]{\sk{अन्यत्र} (anyatra) — elsewhere}

\label{lesson:SA-C95-anyatra}



\begin{quote}
Before the new one: say the Sanskrit for beside, then the Sanskrit for everywhere.
\end{quote}

\subsection*{You'll want to know: \sk{अन्यत्र}}
\textbf{\sk{अन्यत्र}} — \emph{anyatra} — \textquotedblleft{}elsewhere\textquotedblright{}.

\begin{grammarlens}[title={What you've built}]
One more word for this chapter.
\end{grammarlens}

\subsection*{Your turn}
\begin{itemize}
\raggedright
  \item \emph{Say it:} \emph{anyatra}
  \item \emph{Say it:} \emph{anyatra}, once more
  \item \emph{Say it:} say \emph{anyatra}
  \item \emph{Recall:} say the Sanskrit for beside, then the Sanskrit for everywhere, then say \emph{anyatra} again
\end{itemize}

\begin{tcolorbox}[breakable,colback=teal!4,colframe=teal!35!black,title={Before you move on}]
What does \sk{अन्यत्र} mean? (\textquotedblleft{}Elsewhere\textquotedblright{}.) Say it once more.
\end{tcolorbox}
\section[\texorpdfstring{\sk{वीथी} (vīthī)}{वीथी (vīthī)}]{\sk{वीथी} (vīthī) — a street}

\label{lesson:SA-C95-vithi}



\begin{quote}
Before the new one: say the Sanskrit for everywhere, then the Sanskrit for elsewhere.
\end{quote}

\subsection*{You'll want to know: \sk{वीथी}}
\textbf{\sk{वीथी}} — \emph{vīthī} — \textquotedblleft{}a street\textquotedblright{}.

\begin{grammarlens}[title={What you've built}]
One more word for this chapter.
\end{grammarlens}

\subsection*{Your turn}
\begin{itemize}
\raggedright
  \item \emph{Say it:} \emph{vīthī}
  \item \emph{Say it:} \emph{vīthī}, once more
  \item \emph{Say it:} say \emph{vīthī}
  \item \emph{Recall:} say the Sanskrit for everywhere, then the Sanskrit for elsewhere, then say \emph{vīthī} again
\end{itemize}

\begin{tcolorbox}[breakable,colback=teal!4,colframe=teal!35!black,title={Before you move on}]
What does \sk{वीथी} mean? (\textquotedblleft{}A street\textquotedblright{}.) Say it once more.
\end{tcolorbox}
